\addtocontents{toc}{\protect\needspace{5\baselineskip}}
\section{泛化、模型容量与任务适配}
\subsection{训练集、验证集与测试集}

泛化（generalization）是模型在新输入上仍能产生合适结果的能力。训练集用于计算参数更新；验证集用于选择设置和停止训练的位置；测试集用于检查模型对新输入的表现。三者在训练与评价过程中的用途不同。

训练损失持续下降，说明模型在参与参数更新的数据上拟合得更好。图~\ref{l10:fig:training-validation} 同时给出训练损失与验证损失随 epoch 变化的曲线：训练损失继续趋近于零，而验证损失先下降，随后趋于平缓并缓慢上升。这个变化说明，继续改进训练集上的拟合，可能无法同步改进新数据上的结果。

\begin{figure}[H]
\centering
\includegraphics[width=0.84\linewidth]{training_validation.png}
\caption{训练损失与验证损失随训练轮次的变化。纵轴为负对数似然损失；该曲线用于说明泛化现象，与前面数字距离损失的数值例题分开讨论。}
\label{l10:fig:training-validation}
\end{figure}

验证数据为训练停止位置提供依据：参数更新继续使用训练集，停止点与设置则根据验证表现选择。测试数据用于最终检查泛化。由此，优化训练目标、选择模型和评估新输入表现形成相互关联而职责明确的环节。

\subsection{模型容量、欠拟合与过拟合}

\subsubsection{多项式拟合的类比}
模型容量（capacity）描述模型可以表达多复杂的关系。一个简单模型可能连训练样本中的主要变化都无法拟合，表现为欠拟合（underfitting）。增加表达能力后，模型可以捕捉这些变化；若过于贴合训练样本的具体细节，则可能在未见输入上出现较大误差，表现为过拟合（overfitting）。

多项式拟合提供了一个直观例子。过低阶的曲线难以追随数据趋势；合适的阶数能够表达主要关系；过高阶的曲线可以更加贴近已知点，同时在已知点之间或样本范围之外产生较大振荡。这个类比说明，训练误差的改善需要结合新数据表现来判断。

\subsubsection{训练轮次与模型容量是两个观察轴}
图~\ref{l10:fig:capacity} 的横轴是模型容量，纵轴是误差。示意曲线中，容量增加时训练误差下降，而泛化误差在合适容量附近较低，继续增大容量后又上升。两条曲线的差异称为泛化差距（generalization gap）。

\begin{figure}[H]
\centering
\includegraphics[width=0.87\linewidth]{capacity_gap.png}
\caption{模型容量与训练误差、泛化误差之间的示意关系。横轴是容量，不是训练轮次。}
\label{l10:fig:capacity}
\end{figure}

训练轮次衡量同一训练过程推进了多久，模型容量描述结构能够表达多复杂的关系。前一幅曲线用于理解训练停止位置，后一幅曲线用于理解容量选择。两者都需要关注训练数据之外的表现，但对应的可调对象不同。

\subsection{从权重图像理解学习到的特征}

一层直接连接图像像素的网络，其某个单元具有 784 个输入权重。按照像素展开时的顺序，把这些权重重新排列成 $28\times28$ 的网格，就可以观察这个单元如何对不同像素位置加权。这种可视化把权重向量重新映射回输入的空间结构。

\begin{figure}[H]
\centering
\includegraphics[width=0.97\linewidth]{weights_visualization.png}
\caption{单层 MNIST 网络的类别权重图与两层网络的特征权重图。上方以数字 0--9 标注，下方以特征编号标注。}
\label{l10:fig:weights}
\end{figure}

直接面向类别的权重图可以呈现与相应数字外观有关的结构。两层网络的隐藏单元则提供后续输出层要组合的中间特征，因此图中用特征编号而非数字类别为它们命名。这个区别与前面的表示学习一致：输出判断可以建立在多个中间特征的组合上。

\subsection{无免费午餐定理的任务层面含义}

无免费午餐定理（no free lunch theorem）讨论算法在广泛可能问题上的平均表现。一种概括性表述是：当对所有可能的输入生成分布取平均时，各种分类算法在未观测输入上的平均错误率相同。这里的关键限定是\emph{对所有可能分布取平均}。

针对一个具体任务，模型面对的是具有特定结构的数据。网络结构、特征表示和训练设置都需要与这种结构相适配。因而，“从数据中学习”仍然包含任务相关的模型设计与选择；泛化评价用于检验这些选择是否在新输入上有效。平均性能结论依赖所选问题空间及平均方式，并不意味着各算法在某一个具体数据集上的表现相同。
